\documentclass[aspectratio=169]{beamer}

\usepackage[utf8]{inputenc}
\usepackage[T1]{fontenc}
\usepackage{amsmath,amssymb}
\usepackage{listings}
\usepackage{xcolor}
\usepackage{array}
\usepackage{booktabs}
\usepackage{hyperref}

% ── Colour palette ─────────────────────────────────────────────────────────────
\definecolor{orangetheme}{RGB}{220,95,0}
\definecolor{pyblue}    {RGB}{42,115,205}
\definecolor{pygreen}   {RGB}{0,128,0}
\definecolor{pygray}    {RGB}{110,110,110}
\definecolor{pyred}     {RGB}{180,0,0}
\definecolor{pybg}      {RGB}{248,248,252}
\definecolor{linkblue}  {RGB}{0,100,200}
\definecolor{rowred}    {RGB}{255,232,232}
\definecolor{rowblue}   {RGB}{227,238,255}

% ── Beamer theme ───────────────────────────────────────────────────────────────
\usetheme{default}
\setbeamercolor{title}          {fg=orangetheme}
\setbeamercolor{frametitle}     {fg=orangetheme}
\setbeamercolor{structure}      {fg=orangetheme}
\setbeamercolor{itemize item}   {fg=orangetheme}
\setbeamercolor{itemize subitem}{fg=orangetheme}
\setbeamerfont{title}    {size=\Large,series=\bfseries}
\setbeamerfont{frametitle}{size=\large,series=\bfseries}
\setbeamertemplate{navigation symbols}{}
\setbeamertemplate{footline}[frame number]

% ── Python listings style ──────────────────────────────────────────────────────
\lstdefinestyle{python}{
  language        = Python,
  basicstyle      = \footnotesize\ttfamily,
  keywordstyle    = \color{pyblue}\bfseries,
  stringstyle     = \color{pyred},
  commentstyle    = \color{pygreen}\itshape,
  backgroundcolor = \color{pybg},
  frame           = single,
  framerule       = 0.4pt,
  rulecolor       = \color{gray!40},
  showstringspaces= false,
  breaklines      = true,
  xleftmargin     = 4pt,
  xrightmargin    = 4pt,
  aboveskip       = 4pt,
  belowskip       = 4pt,
  morekeywords    = {True,False,None,sat,unsat},
}

% ── Meta ───────────────────────────────────────────────────────────────────────
\title{Rules of Inference}
\subtitle{Discrete Mathematics --- Lesson 5}
\author{Notes By Pr. El Hadiq Zouhair}
\date{}

% ══════════════════════════════════════════════════════════════════════════════
\begin{document}

\begin{frame}
  \titlepage
\end{frame}

% ── 1. Overview ────────────────────────────────────────────────────────────────
\begin{frame}{Overview}
  \begin{itemize}
    \item Most of the time, the full explanation is not needed.
    \item Inference is the process of deriving logical conclusions from known premises.
    \item Aristotle classically divided it in two: \textbf{deduction} and \textbf{induction}.
  \end{itemize}
\end{frame}

% ── 2. Motivation ──────────────────────────────────────────────────────────────
\begin{frame}[fragile]{Motivation}
  Assume the following statements as true:
  \begin{itemize}
    \item If you have cherries, you can eat. \quad $\mathit{have} \Rightarrow \mathit{eat}$
    \item You have cherries. \quad $\mathit{have}$
  \end{itemize}
  Intuitively, by inference, this implies you can eat. But how do you get there?

  \begin{lstlisting}[style=python]
from sympy import symbols, And, Implies, simplify_logic
from sympy.logic.boolalg import truth_table

have, eat = symbols('have eat')
expr = And(Implies(have, eat), have)
print(simplify_logic(expr))          # eat & have

for vals, result in truth_table(expr, [have, eat]):
    print(list(vals), '->', result)
  \end{lstlisting}

  \smallskip
  \begin{center}
    \renewcommand{\arraystretch}{1.3}
    \begin{tabular}{|c|c|c|}
      \hline
      \textbf{have} & \textbf{eat} & \textbf{eat $\wedge$ have} \\
      \hline
      \rowcolor{rowred}  True  & True  & True  \\
      \rowcolor{rowblue} True  & False & False \\
      \rowcolor{rowblue} False & True  & False \\
      \rowcolor{rowblue} False & False & False \\
      \hline
    \end{tabular}
  \end{center}
\end{frame}

% ── 3. Limits of Truth Tables ──────────────────────────────────────────────────
\begin{frame}[fragile]{Limits of Truth Tables}
  Simplification through equivalences is often insufficient to obtain conclusive results:

  \begin{lstlisting}[style=python]
from sympy import symbols, Or, And, simplify_logic
from sympy.logic.boolalg import truth_table
from itertools import combinations

a, b, c, d, e = symbols('a b c d e')
# BooleanCountingFunction[2,5]: at least 2 of 5 variables are True
expr = Or(*[And(x, y) for x, y in combinations([a,b,c,d,e], 2)])
simplified = simplify_logic(expr)
print(simplified)   # still complex

# Truth table has 2^5 = 32 rows -- too large to inspect by hand
for vals, result in truth_table(simplified, [a, b, c, d, e]):
    print(list(vals), '->', result)
  \end{lstlisting}

  \medskip
  With 5 variables, a truth table has $2^5 = 32$ rows.
  Rules of inference provide a structured shortcut.
\end{frame}

% ── 4. Rules of Inference ──────────────────────────────────────────────────────
\begin{frame}[fragile]{Rules of Inference}
  The most common rules are:

  \medskip
  \begin{center}
    \renewcommand{\arraystretch}{1.35}
    \small
    \begin{tabular}{|l|c|c|c|}
      \hline
      \textbf{Rule} & \textbf{Premises} & & \textbf{Conclusion} \\
      \hline
      Modus Ponens           & $p \;\wedge\; (p\to q)$      & $\to$ & $q$       \\
      Modus Tollens          & $\neg q \;\wedge\; (p\to q)$ & $\to$ & $\neg p$  \\
      Hypothetical Syllogism & $(p\to q) \;\wedge\; (q\to r)$& $\to$ & $p\to r$ \\
      Disjunctive Syllogism  & $(p\vee q) \;\wedge\; \neg p$ & $\to$ & $q$      \\
      Addition               & $p$                           & $\to$ & $p\vee q$ \\
      Simplification         & $p\wedge q$                   & $\to$ & $p$       \\
      Conjunction            & $p \;\wedge\; q$              & $\to$ & $p\wedge q$ \\
      Resolution             & $(p\vee q) \;\wedge\; (\neg p\vee r)$ & $\to$ & $q\vee r$ \\
      \hline
    \end{tabular}
  \end{center}

  \medskip
  \begin{lstlisting}[style=python]
from sympy import symbols, And, Or, Not, Implies
from sympy.logic.inference import satisfiable

p, q, r = symbols('p q r')
# Verify: premises & ~conclusion must be UNSAT for a valid rule
checks = {
    'Modus Ponens':     And(p, Implies(p,q), Not(q)),
    'Modus Tollens':    And(Not(q), Implies(p,q), p),
    'Resolution':       And(Or(p,q), Or(Not(p),r), Not(Or(q,r))),
}
for name, f in checks.items():
    print(name, ':', not satisfiable(f))   # True => valid
  \end{lstlisting}
\end{frame}

% ── 5. Example 1: Swimming ─────────────────────────────────────────────────────
\begin{frame}[fragile]{Example 1}
  Assume the following statements as true:
  \begin{itemize}
    \item You will go swim \textbf{or} go to the gym. \quad $\mathrm{Swim} \vee \mathrm{Gym}$
    \item If you go swimming, you will need a swimsuit. \quad $\mathrm{Swim} \to \mathrm{Swimsuit}$
  \end{itemize}

  Is it true that you will need a swimsuit \textbf{or} go to the gym?
  $\mathrm{Swimsuit} \vee \mathrm{Gym}$

  \medskip
  Using a conditional equivalence: $\neg\,\mathrm{Swim} \vee \mathrm{Swimsuit}$ \\
  Using \textbf{resolution}: $\textcolor{orangetheme}{\mathrm{Swimsuit} \vee \mathrm{Gym}}$

  \begin{lstlisting}[style=python]
from sympy import symbols, And, Or, Not, Implies
from sympy.logic.inference import satisfiable

Swim, Gym, Swimsuit = symbols('Swim Gym Swimsuit')
premises = And(Or(Swim, Gym), Implies(Swim, Swimsuit))
conclusion = Or(Swimsuit, Gym)

# Valid if premises & ~conclusion is unsatisfiable
print(not satisfiable(And(premises, Not(conclusion))))  # True
  \end{lstlisting}
\end{frame}

% ── 6. Example 2: NFT Setup ────────────────────────────────────────────────────
\begin{frame}{Example 2}
  Assume the following statements as true:
  \begin{itemize}
    \item Your Invisible Rock \#14 NFT can be sold for $\approx$4.32 ETH. \\
          $\mathrm{Sell(Rock)} \to \mathrm{Get(4.32\,ETH)}$
    \item 4.32 ETH can be sold for \$7000. \quad
          $\mathrm{Get(4.32\,ETH)} \to \mathrm{Get(\$7000)}$
    \item If you have \$10\,000, you can buy a car. \quad
          $\mathrm{Have(\$10\,000)} \to \mathrm{Buy(Car)}$
    \item You saved \$3000 for a car. \quad $\mathrm{Savings(\$3000)}$
  \end{itemize}

  \medskip
  If you sell your Invisible Rock ($\mathrm{Sell(Rock)}$), can you buy yourself a car?
\end{frame}

% ── 7. Example 2: NFT Solution ─────────────────────────────────────────────────
\begin{frame}[fragile]{Example 2 (continued)}
  \begin{enumerate}
    \item \textbf{Hypothetical syllogism}: \quad $\mathrm{Sell(Rock)} \to \mathrm{Get(\$7000)}$
    \item \textbf{Modus ponens}: \quad $\mathrm{Get(\$7000)}$
    \item \textbf{Conjunction}: \quad $\mathrm{Get(\$7000)} \wedge \mathrm{Savings(\$3000)}$
    \item \textbf{Addition + Modus ponens}: \quad $\textcolor{orangetheme}{\mathrm{Have(\$10\,000)}}$
    \item \textbf{Modus ponens}: \quad $\textcolor{orangetheme}{\mathrm{Buy(Car)}}$
  \end{enumerate}

  \begin{lstlisting}[style=python]
from sympy import symbols, And, Implies, Not
from sympy.logic.inference import satisfiable

SellRock, Get7000, Have10000, BuyCar, Savings = symbols(
    'SellRock Get7000 Have10000 BuyCar Savings')

premises = And(
    SellRock,
    Implies(SellRock, Get7000),                  # hyp. syllogism (steps 1&2)
    Savings,
    Implies(And(Get7000, Savings), Have10000),    # addition principle
    Implies(Have10000, BuyCar),
)
print(not satisfiable(And(premises, Not(BuyCar))))  # True
  \end{lstlisting}
\end{frame}

% ── 8. Rules for Quantified Statements ─────────────────────────────────────────
\begin{frame}[fragile]{Rules for Quantified Statements}
  Inference is also used for quantifiers:

  \medskip
  \begin{center}
    \renewcommand{\arraystretch}{1.35}
    \small
    \begin{tabular}{|l|c|c|c|}
      \hline
      \textbf{Rule} & \textbf{Premise} & & \textbf{Conclusion} \\
      \hline
      Universal Instantiation   & $\forall x\,P(x)$        & $\to$ & $P(c)$ \\
      Universal Generalization  & $P(c)$ for arbitrary $c$ & $\to$ & $\forall x\,P(x)$ \\
      Existential Instantiation & $\exists x\,P(x)$         & $\to$ & $P(c)$ for some $c$ \\
      Existential Generalization& $P(c)$ for some $c$      & $\to$ & $\exists x\,P(x)$ \\
      \hline
    \end{tabular}
  \end{center}

  All humans live on Earth: $\forall x\,(H(x) \to L(x))$. \quad You are human: $H(\text{you})$.

  \begin{lstlisting}[style=python]
from z3 import DeclareSort, Function, BoolSort, Const, ForAll, Implies
from z3 import Solver, Not

Person = DeclareSort('Person')
you = Const('you', Person)
x   = Const('x',   Person)
H = Function('H', Person, BoolSort())
L = Function('L', Person, BoolSort())

s = Solver()
s.add(ForAll([x], Implies(H(x), L(x))))  # all x: H(x)->L(x)
s.add(H(you))                             # H(you)
s.add(Not(L(you)))                        # assume ~L(you)
print(s.check())   # unsat  =>  L(you) is derived
  \end{lstlisting}
\end{frame}

% ── 9. Example: Exam & Director ────────────────────────────────────────────────
\begin{frame}[fragile]{Example: Exam \& Director}
  Assume the following statements as true:
  \begin{itemize}
    \item If no one cheated, no one meets the director. \quad
          $\neg\exists x\,\mathrm{Cheat}(x) \to \neg\exists x\,\mathrm{Director}(x)$
    \item Octave (class president) met the director. \quad
          $\mathrm{Director(Octave)}$ \quad
          $\xrightarrow{\text{ex. gen.}}$ \; $\exists x\,\mathrm{Director}(x)$
    \item No one cheated or someone gets detention. \quad
          $\neg\exists x\,\mathrm{Cheat}(x) \;\vee\; \exists x\,\mathrm{Detention}(x)$
  \end{itemize}

  Will someone get a detention?

  \begin{lstlisting}[style=python]
from sympy import symbols, And, Or, Not, Implies
from sympy.logic.inference import satisfiable

ExCheat, ExDirector, ExDetention = symbols(
    'ExCheat ExDirector ExDetention')

premises = And(
    Implies(Not(ExCheat), Not(ExDirector)),   # ~ExCheat -> ~ExDirector
    ExDirector,                               # Director(Octave) => ExDirector
    Or(Not(ExCheat), ExDetention),            # ~ExCheat v ExDetention
)
# Modus tollens: ExDirector => ExCheat; disj. syllogism => ExDetention
print(not satisfiable(And(premises, Not(ExDetention))))  # True
  \end{lstlisting}
\end{frame}

% ── 10. Summary ────────────────────────────────────────────────────────────────
\begin{frame}{Summary}
  \begin{itemize}
    \item The main rules of inference are:
    \begin{itemize}
      \item \textbf{Modus ponens}, \textbf{modus tollens}, \textbf{hypothetical syllogism},
            \textbf{disjunctive syllogism}, \textbf{addition}, \textbf{simplification},
            \textbf{conjunction} and \textbf{resolution}.
    \end{itemize}
    \item Quantified rules of inference:
    \begin{itemize}
      \item \textbf{Universal instantiation}, \textbf{universal generalization},
            \textbf{existential instantiation} and \textbf{existential generalization}.
    \end{itemize}
    \item These rules make it easy to extract useful conclusions out of compound propositions.
    \item Verify rules with \texttt{sympy.logic.inference.satisfiable};
          quantified rules with \texttt{z3.Solver}.
    \item The next lesson goes beyond true/false values --- onward with \textbf{sets}.
  \end{itemize}
\end{frame}

\end{document}
